% !TEX root = ../analisi-dei-capitolati.tex
% Capitolo 6 - Altri capitolati (C3-C9)

\section{Altri capitolati}\label{sec:altri}

I capitolati C3--C9 sono stati letti e discussi da tutto il gruppo nella
fase di prima selezione, ma non sono entrati tra i finalisti e non sono
stati approfonditi con i proponenti. Per ciascuno si riportano obiettivo,
requisiti e tecnologie principali e la valutazione del gruppo.

% ---------------------------------------------------------------------
\subsection{C3 -- FossDot}
\textbf{Proponente}: Uranion EMEA Ltd.

\textbf{Obiettivo.} Gli smart speaker in commercio sono dispositivi chiusi,
legati al cloud del produttore e a un solo assistente vocale. Uranion
propone un ecosistema open source completo per un assistente vocale
domestico sul modello dell'Echo Dot, in cui il proprietario sceglie il
provider di intelligenza artificiale (in cloud o locale), collega servizi
esterni tramite \glo{MCP} e fa cooperare più dispositivi in casa.

\textbf{Requisiti e tecnologie.} L'MVP comprende la scelta e la validazione
dell'hardware di riferimento (scheda, array di microfoni, altoparlante) con
distinta base e prototipo, un firmware che cattura l'audio, riconosce la
parola di attivazione e trasmette la voce al server, un servizio di
\emph{bridging} con pipeline speech-to-text, LLM e text-to-speech, e
un'interfaccia web per registrare i dispositivi. La consegna completa
aggiunge più provider commutabili, integrazione MCP, \emph{skill} e
agenti, funzioni multi-dispositivo e gestione sicura dei segreti, con una
copertura dei test software dell'80\%. Lo stack è libero, con licenza
Apache~2.0 o MIT; Rust è preferito per il firmware, e sono suggeriti
Raspberry Pi, ReSpeaker, openWakeWord, faster-whisper, Piper e Ollama.

\textbf{Valutazione.} Il tema della privacy e del controllo sui propri dati
è stato apprezzato. Il progetto richiede però competenze di elettronica e
di programmazione embedded lontane dal percorso del gruppo, l'acquisto e la
gestione di hardware e una dimostrazione finale su dispositivo fisico: un
rischio ritenuto non sostenibile.

% ---------------------------------------------------------------------
\subsection{C4 -- Piattaforma per la gestione dei veicoli e dei servizi post-vendita}
\textbf{Proponente}: M31 S.r.l.

\textbf{Obiettivo.} Le informazioni su modelli, veicoli venduti,
manutenzioni e campagne di richiamo sono sparse tra sistemi diversi. M31
chiede una piattaforma distribuita che colleghi casa automobilistica,
concessionarie e proprietari, con accesso separato ai dati in base al ruolo
e all'organizzazione.

\textbf{Requisiti e tecnologie.} Sono obbligatori la gestione di modelli,
veicoli, concessionarie e proprietari, le richieste di manutenzione con
aggiornamento dello stato, le campagne di richiamo con individuazione dei
veicoli coinvolti, le notifiche, tre interfacce web distinte e \glo{API}
sicure, documentate e versionate; i dati dei veicoli possono essere
inseriti a mano o simulati. Tra gli opzionali figurano storico tecnico,
passaggio di proprietà, preventivi, manutenzione predittiva e ricerca
geografica. Sono richiesti anche un'analisi dello stato dell'arte e un
\emph{test book}. Le tecnologie indicate sono microservizi in Node.js con
Nest.js e TypeScript, Go per le componenti più esigenti, NATS o Kafka per
gli eventi, Kubernetes su Google Cloud Platform, PostgreSQL, MongoDB,
Redis, Angular, Prometheus e Grafana.

\textbf{Valutazione.} Il capitolato è chiaro e ben strutturato, ma la parte
applicativa è un gestionale tradizionale e la complessità sta soprattutto
nell'infrastruttura distribuita (microservizi, code di eventi,
orchestrazione). Il gruppo ha preferito un progetto in cui l'innovazione
riguardi il comportamento del prodotto.

% ---------------------------------------------------------------------
\subsection{C5 -- RareCare}
\textbf{Proponente}: Miriade S.r.l.

\textbf{Obiettivo.} Le famiglie di persone con malattie rare devono
coordinarsi da sole con centri clinici, pediatri, terapisti, scuola e
servizi sociali. RareCare Navigator, pensato inizialmente per la Sindrome
di Sanfilippo, vuole affiancare la figura del \emph{Care Navigator} con un
assistente basato sull'intelligenza artificiale, un sistema di scadenze e
un archivio condiviso dei documenti di cura.

\textbf{Requisiti e tecnologie.} Sono obbligatori un modello di ruoli e
permessi granulare, un motore di scadenze e notifiche (rinnovi, piani
assistenziali ed educativi) e un'interfaccia conversazionale \glo{RAG} che
risponde solo sulla base di fonti ufficiali verificate; sono opzionali un
assistente vocale, la mappatura dei servizi territoriali e l'analisi
predittiva dei bisogni. Il nucleo della piattaforma deve essere
indipendente dalla patologia, con pacchetti di configurazione per
estenderla ad altre malattie. Sono raccomandati un'architettura serverless
su AWS (API Gateway, Lambda, Bedrock), un'app multipiattaforma in Flutter o
React Native, cifratura dei dati, autenticazione a più fattori e
accessibilità WCAG~2.1~AA.

\textbf{Valutazione.} Il valore sociale del progetto è stato riconosciuto
da tutto il gruppo e la componente RAG è affine a quella di C2. Hanno
pesato però la delicatezza del dominio (dati sanitari di minori,
conformità GDPR, possibili implicazioni regolatorie), l'ampiezza delle
responsabilità richieste e uno sviluppo mobile multipiattaforma che il
gruppo non intendeva affrontare.

% ---------------------------------------------------------------------
\subsection{C6 -- Caching semantico e dashboard FinOps per LiteLLM}
\textbf{Proponente}: Var Group S.p.A.

\textbf{Obiettivo.} Var Group usa il progetto open source LiteLLM come
orchestratore davanti a più provider di modelli linguistici. In produzione
ne ha riscontrato due limiti: la cache riconosce solo richieste identiche,
quindi domande equivalenti formulate in modo diverso vengono pagate due
volte, e non esiste una visione dei costi per team, progetto o utente.

\textbf{Requisiti e tecnologie.} Il progetto chiede di estendere LiteLLM
con una cache semantica basata su \glo{embedding} e ricerca per
similarità, con soglia configurabile e invalidazione, e con un servizio
FinOps che attribuisca i costi, gestisca budget e avvisi e preveda la
spesa mensile, con isolamento dei dati tra team. Prima dello sviluppo sono
richiesti una sessione di design thinking, lo user story mapping e un PoC
di fattibilità; i test devono coprire almeno il 70\% del codice. Le
tecnologie consigliate sono Python con FastAPI, Docker Compose,
sentence-transformers o gli embedding di OpenAI, Redis, Qdrant o
pgvector, PostgreSQL e React con Chart.js.

\textbf{Valutazione.} Il tema è attuale e ben delimitato, ma molto
orientato all'infrastruttura: il prodotto non ha un utente finale diretto
e il suo valore si misura soprattutto in risparmi di costo. L'interesse
del gruppo è stato inferiore rispetto ai finalisti.

% ---------------------------------------------------------------------
\subsection{C7 -- SmartDashboard}
\textbf{Proponente}: Ergon Informatica S.r.l.

\textbf{Obiettivo.} I dati delle aziende di produzione e distribuzione
restano spesso chiusi nel gestionale e ogni nuovo report richiede
l'intervento del reparto IT. Ergon chiede una piattaforma di
\glo{Business Intelligence} self-service che recuperi i dati dal
gestionale, li consolidi in un \glo{data warehouse} e permetta agli utenti
di costruire da soli le proprie dashboard.

\textbf{Requisiti e tecnologie.} La web app deve offrire un costruttore di
dashboard a trascinamento con grafici, KPI, tabelle pivot, gauge e mappe,
filtri trasversali, esportazione in PDF ed Excel e gestione degli utenti.
Un livello di intelligenza artificiale deve generare visualizzazioni da
richieste in linguaggio naturale e, come estensioni progressive, produrre
insight automatici, suggerire widget e spiegare gli indicatori. Il
proponente consiglia di costruire prima la pipeline di Business
Intelligence e poi il livello AI, e può fornire i dati di un caso di studio
reale. Le tecnologie suggerite sono DuckDB per il data warehouse (con
supporto desiderabile ad altri database), PostgreSQL o MariaDB, Polars per
le trasformazioni, React, Angular o Vue e un LLM commerciale.

\textbf{Valutazione.} Il capitolato è chiaro e propone un percorso
incrementale ragionevole. La parte prevalente del lavoro riguarda però la
pipeline dei dati e il costruttore di dashboard, ambiti che hanno
suscitato meno interesse rispetto agli assistenti conversazionali dei
finalisti.

% ---------------------------------------------------------------------
\subsection{C8 -- Sonar}
\textbf{Proponente}: Halue S.r.l. Società Benefit.

\textbf{Obiettivo.} Con gli anni uno store Shopify accumula app, webhook,
campi personalizzati, script e codice di tema di cui nessuno ricorda più
lo scopo, e questa conoscenza si perde a ogni cambio di fornitore. Sonar
deve collegarsi a uno store, acquisirne configurazioni e codice e
ricostruire in automatico cosa è installato, cosa reagisce a quali
eventi, quali campi sono davvero usati e cosa è orfano.

\textbf{Requisiti e tecnologie.} Le scansioni sono lunghe e soggette ai
limiti di frequenza delle API di Shopify: devono essere durevoli,
riprendibili e osservabili, e salvare ogni volta uno snapshot immutabile.
I dati confluiscono in un grafo di entità e dipendenze, su cui un motore di
regole estensibile segnala codice morto e configurazioni orfane con una
motivazione per ciascuna segnalazione; la documentazione è esportabile in
Markdown o HTML. Sono richiesti cifratura dei token, isolamento tra
organizzazioni, copertura dei test del 70\% e un \emph{test book}. Lo
stack è in gran parte vincolante: Next.js su Vercel, Vercel Workflows e
Cron Jobs, Supabase con autenticazione Google, TypeScript e Shopify Admin
GraphQL API.

\textbf{Valutazione.} È tra i capitolati più curati per requisiti e criteri
di accettazione, ma è legato a un ecosistema specifico con molti vincoli
tecnologici, e il dominio dell'e-commerce enterprise non ha coinvolto il
gruppo quanto i finalisti.

% ---------------------------------------------------------------------
\subsection{C9 -- SpecCheck}
\textbf{Proponente}: Zucchetti S.p.A.

\textbf{Obiettivo.} Con lo sviluppo assistito dagli LLM il collo di
bottiglia si è spostato dalla scrittura del codice alla sua verifica:
compilatori e test dicono se il codice funziona, non se fa ciò che era
stato chiesto. Zucchetti chiede di estendere l'agente di sviluppo
OpenCode con strumenti che verifichino la tracciabilità tra casi d'uso,
requisiti e codice sorgente.

\textbf{Requisiti e tecnologie.} Il prodotto è un pacchetto di skill,
comandi e almeno un plugin per OpenCode, installabile con uno script, che
legga i diagrammi e le descrizioni dei casi d'uso (per esempio in Mermaid,
XMI o Markdown) e le specifiche dei requisiti, li confronti tra loro e con
il comportamento del codice e mostri le discrepanze in una dashboard,
anche testuale. Il confronto finale, che richiede precisione, va
realizzato come plugin o programma esterno e non affidato solo al modello.
Sono opzionali una dashboard grafica, l'individuazione del codice non
coperto da specifiche e la lettura diretta da PDF. Il proponente fornisce
accesso a modelli di grandi dimensioni, con un tetto di spesa, e a
modelli più piccoli senza limiti.

\textbf{Valutazione.} Il tema è originale e molto vicino agli argomenti del
corso, ma il capitolato si dichiara fortemente esplorativo: parte dei
risultati dipende da capacità degli LLM non ancora consolidate. Il gruppo
ha preferito un progetto con perimetro e criteri di accettazione più
definiti.
